\documentclass[11pt]{article}
\usepackage[margin=1in]{geometry}
\usepackage{amsmath}
\usepackage{amssymb}
\usepackage{amsthm}
\usepackage{booktabs}
\usepackage{longtable}
\usepackage{array}
\usepackage{hyperref}
\usepackage{xcolor}
\usepackage{listings}
\usepackage{enumitem}
\usepackage{fancyhdr}
\usepackage{titlesec}

\hypersetup{
    colorlinks=true,
    linkcolor=blue,
    urlcolor=blue,
    citecolor=blue
}

\lstset{
    language=R,
    basicstyle=\ttfamily\small,
    keywordstyle=\color{blue},
    commentstyle=\color{gray},
    stringstyle=\color{red},
    breaklines=true,
    frame=single,
    backgroundcolor=\color{gray!10}
}

\theoremstyle{definition}
\newtheorem{definition}{Definition}
\newtheorem{theorem}{Theorem}
\newtheorem{proposition}{Proposition}
\newtheorem{remark}{Remark}

\title{Mathematical Foundations of N-of-1 Trial Simulation:\\
A Comprehensive Analysis}
\author{pmsimstats2025 Project Documentation}
\date{December 17, 2025}

\begin{document}

\maketitle

\begin{abstract}
This document provides an in-depth, comprehensive synthesis of all mathematical
elements underlying the \texttt{pmsimstats2025} N-of-1 clinical trial simulation
framework. The analysis draws from extensive documentation and represents a
unified treatment of the project's theoretical foundations, covering covariance
matrix architecture, positive definiteness constraints, correlation structures,
carryover modeling, and mixed-effects model specification.
\end{abstract}

\tableofcontents
\newpage

%==============================================================================
\part{Covariance Matrix Architecture}
%==============================================================================

\section{The Core Identity: $\Sigma = DRD$}

The fundamental mathematical identity governing the simulation is the
decomposition of the covariance matrix into correlation and variance components:

\begin{equation}
\boxed{\Sigma = D \cdot R \cdot D}
\end{equation}

where:
\begin{itemize}
    \item $\Sigma$ is the full covariance matrix ($26 \times 26$ for 8-timepoint designs)
    \item $R$ is the correlation matrix (dimensionless, all diagonal entries $= 1$)
    \item $D$ is a diagonal matrix of standard deviations
\end{itemize}

This decomposition separates two distinct concerns:
\begin{enumerate}
    \item \textbf{Correlation structure ($R$)}: How variables relate to each other
    (independent of scale)
    \item \textbf{Variance magnitude ($D$)}: The scale of each variable's variability
\end{enumerate}

\section{Block Partitioning Strategy}

The full covariance matrix is partitioned into blocks for computational
efficiency:

\begin{equation}
\Sigma = \begin{bmatrix}
\Sigma_{11} & \Sigma_{12} \\
\Sigma_{12}' & \Sigma_{22}
\end{bmatrix}
\end{equation}

\begin{table}[h]
\centering
\caption{Block Matrix Components}
\begin{tabular}{@{}lll@{}}
\toprule
\textbf{Block} & \textbf{Dimension} & \textbf{Contents} \\
\midrule
$\Sigma_{11}$ & $(3n_{tp}) \times (3n_{tp})$ & Response components (BR, ER, TR) at all timepoints \\
$\Sigma_{22}$ & $2 \times 2$ & Biomarker and baseline \\
$\Sigma_{12}$ & $(3n_{tp}) \times 2$ & Cross-covariances \\
\bottomrule
\end{tabular}
\end{table}

For 8-timepoint designs: $\Sigma_{11}$ is $24 \times 24$, giving a full matrix
of $26 \times 26$.

\section{Intra-Block Structure}

Within $\Sigma_{11}$, there is further structure:

\begin{equation}
\Sigma_{11} = \begin{bmatrix}
\Sigma_{BR} & C_{BR,ER} & C_{BR,TR} \\
C_{ER,BR} & \Sigma_{ER} & C_{ER,TR} \\
C_{TR,BR} & C_{TR,ER} & \Sigma_{TR}
\end{bmatrix}
\end{equation}

Each diagonal block (e.g., $\Sigma_{BR}$) is an $n_{tp} \times n_{tp}$ matrix with:
\begin{itemize}
    \item Diagonal entries $= 1$ (unit variance for correlation matrix)
    \item Off-diagonal entries $=$ autocorrelation parameter (e.g., $c_{br} = 0.75$)
\end{itemize}

Cross-blocks (e.g., $C_{BR,ER}$) contain:
\begin{itemize}
    \item Diagonal entries $= c_{cf1t}$ (same-timepoint cross-correlation)
    \item Off-diagonal entries $= c_{cfct}$ (different-timepoint cross-correlation)
\end{itemize}

%==============================================================================
\part{Conditional Multivariate Normal Sampling}
%==============================================================================

\section{The Two-Stage Algorithm}

Rather than inverting the full $26 \times 26$ matrix, the simulation uses the
Schur complement to enable efficient two-stage sampling:

\subsection{Stage 1: Generate Baseline Variables}

Generate (biomarker, baseline) from $\Sigma_{22}$:
\begin{equation}
(x_2) \sim \mathcal{N}(\mu_2, \Sigma_{22})
\end{equation}

\subsection{Stage 2: Generate Response Components}

Generate response components conditional on Stage 1:
\begin{equation}
x_1 | x_2 \sim \mathcal{N}(\mu_{1|2}, \Sigma_{cond})
\end{equation}

where:
\begin{align}
\mu_{1|2} &= \mu_1 + \Sigma_{12} \Sigma_{22}^{-1} (x_2 - \mu_2) \\
\Sigma_{cond} &= \Sigma_{11} - \Sigma_{12} \Sigma_{22}^{-1} \Sigma_{12}'
\end{align}

\section{Computational Advantages}

This approach reduces computational complexity significantly:

\begin{table}[h]
\centering
\caption{Computational Comparison}
\begin{tabular}{@{}lll@{}}
\toprule
\textbf{Operation} & \textbf{Full Matrix} & \textbf{Two-Stage} \\
\midrule
Matrix inversion & $26 \times 26$ & $2 \times 2$ \\
Cholesky decomposition & $O(26^3)$ & $O(2^3) + O(24^3)$ \\
Numerical stability & Lower & Higher \\
\bottomrule
\end{tabular}
\end{table}

The key insight is that $\Sigma_{22}$ is only $2 \times 2$, making its
inversion trivial and numerically stable.

\section{Conditional Covariance Requirements}

The conditional covariance $\Sigma_{cond}$ is the central matrix for data
generation. It must be:
\begin{enumerate}
    \item \textbf{Positive definite} (all eigenvalues $> 0$)
    \item \textbf{Well-conditioned} (condition number $\kappa < 100$)
    \item \textbf{Numerically stable} for Cholesky decomposition
\end{enumerate}

%==============================================================================
\part{Positive Definiteness Constraints}
%==============================================================================

\section{Mathematical Definition}

\begin{definition}[Positive Definiteness]
A symmetric matrix $\Sigma$ is positive definite if and only if:
\begin{equation}
\mathbf{v}^T \Sigma \mathbf{v} > 0 \quad \forall \mathbf{v} \neq \mathbf{0}
\end{equation}
Equivalently: All eigenvalues $\lambda_i > 0$.
\end{definition}

\section{Verification Methods}

Three methods are used to verify positive definiteness:

\subsection{Method 1: Sylvester's Criterion}

\begin{theorem}[Sylvester's Criterion]
A symmetric matrix is positive definite if and only if all leading principal
minors are positive:
\begin{equation}
\det(\Sigma_k) > 0 \quad \text{for } k = 1, 2, \ldots, n
\end{equation}
where $\Sigma_k$ is the $k \times k$ upper-left submatrix.
\end{theorem}

\subsection{Method 2: Eigenvalue Criterion}

\begin{equation}
\lambda_{min}(\Sigma) > 0
\end{equation}

This is the primary method used in \texttt{build\_sigma\_matrix()}.

\subsection{Method 3: Cholesky Decomposition}

If $\Sigma = LL^T$ exists (with $L$ lower triangular, positive diagonal),
then $\Sigma$ is PD.

\section{Gershgorin Circle Theorem}

\begin{theorem}[Gershgorin]
For eigenvalue bounds:
\begin{equation}
\lambda_i \in \bigcup_{j=1}^{n} \left\{ z \in \mathbb{C} : |z - a_{jj}| \leq R_j \right\}
\end{equation}
where $R_j = \sum_{k \neq j} |a_{jk}|$ is the off-diagonal row sum.
\end{theorem}

For correlation matrices (diagonal $= 1$):
\begin{equation}
\lambda_i \in [1 - R_j, 1 + R_j]
\end{equation}

This implies that if off-diagonal sums exceed 1, negative eigenvalues become
possible.

\section{The Correlation Hierarchy Constraint}

The critical constraint for PD preservation is the correlation hierarchy:

\begin{equation}
\boxed{c_{cfct} < c_{cf1t} < \min(c_{tv}, c_{pb}, c_{br})}
\end{equation}

\begin{remark}
Measurements at different times cannot be more correlated than measurements
at the same time. Violating this creates impossible constraint systems.
\end{remark}

\textbf{Hendrickson's values satisfy this}:
\begin{itemize}
    \item $c_{cfct} = 0.05 < c_{cf1t} = 0.12 < c_{br} = 0.75$ \checkmark
\end{itemize}

\section{Biomarker Correlation Bound}

Empirical analysis establishes:
\begin{equation}
c_{bm} \leq 0.6
\end{equation}
for designs with 8 timepoints. Higher values risk non-PD matrices.

%==============================================================================
\part{Condition Number Analysis}
%==============================================================================

\section{Definition and Interpretation}

The condition number measures numerical stability:

\begin{equation}
\kappa = \frac{\lambda_{max}}{\lambda_{min}}
\end{equation}

\begin{table}[h]
\centering
\caption{Condition Number Interpretation}
\begin{tabular}{@{}lll@{}}
\toprule
\textbf{$\kappa$ Value} & \textbf{Status} & \textbf{Interpretation} \\
\midrule
$\kappa < 10$ & Excellent & Near-perfect numerical stability \\
$10 < \kappa < 100$ & Good & Acceptable for most applications \\
$100 < \kappa < 1000$ & Poor & Numerical precision concerns \\
$\kappa > 1000$ & Severe & Cholesky decomposition unreliable \\
\bottomrule
\end{tabular}
\end{table}

\section{The AR(1) Problem}

AR(1) correlation structure: $\rho_{ij} = \rho^{|t_i - t_j|}$ with $\rho = 0.75$

For evenly-spaced designs (weeks 2, 8, 14, 20), time gaps are 6, 12, 18 weeks:

\begin{table}[h]
\centering
\caption{AR(1) Correlation Decay}
\begin{tabular}{@{}ll@{}}
\toprule
\textbf{Time Gap} & \textbf{AR(1) Correlation} \\
\midrule
6 weeks & $0.75^6 \approx 0.178$ \\
12 weeks & $0.75^{12} \approx 0.032$ \\
18 weeks & $0.75^{18} \approx 0.0056$ \\
\bottomrule
\end{tabular}
\end{table}

\textbf{Correlation ratio}: $0.178 / 0.0056 \approx 32:1$

This extreme ratio creates condition numbers $\kappa \approx 600$--$14,621$ for
evenly-spaced designs.

\section{Why Clustering Solves This}

Clustered designs (weeks 4, 8, 9, 10, 11, 12, 16, 20) have mixed gaps:
\begin{itemize}
    \item Short gaps: 1--4 weeks $\rightarrow$ correlations 0.32--0.75
    \item Long gaps: 8--16 weeks $\rightarrow$ correlations 0.10--0.32
\end{itemize}

This creates \textbf{intermediate correlation values} spanning the correlation
space more evenly, producing balanced eigenvalue distributions with
$\kappa \approx 56$--$61$.

\section{Alternative Correlation Structures}

Five structures are implemented for flexible design:

\begin{table}[h]
\centering
\caption{Correlation Structure Comparison}
\begin{tabular}{@{}lll@{}}
\toprule
\textbf{Structure} & \textbf{Formula} & \textbf{Best For} \\
\midrule
AR(1) & $\rho^{|t_i - t_j|}$ & Clustered designs \\
Exponential & $\exp(-\lambda|t_i - t_j|)$ & Evenly-spaced ($\kappa \approx 18$) \\
Power Law & $(1 + \alpha|t_i - t_j|)^{-\beta}$ & Evenly-spaced ($\kappa \approx 13$) \\
Mat\'ern & $(1 + \sqrt{3}d/\rho)\exp(-\sqrt{3}d/\rho)$ & Flexible smoothness \\
Rational Quadratic & $(1 + t^2/(2\alpha\lambda^2))^{-\alpha}$ & Best stability ($\kappa \approx 10$) \\
\bottomrule
\end{tabular}
\end{table}

%==============================================================================
\part{Three-Factor Response Model}
%==============================================================================

\section{Response Decomposition}

Each participant's response is decomposed into three components:

\begin{equation}
\text{Response}_t = \text{Baseline} + \text{BR}_t + \text{ER}_t + \text{TR}_t
\end{equation}

\begin{table}[h]
\centering
\caption{Response Component Definitions}
\begin{tabular}{@{}llll@{}}
\toprule
\textbf{Component} & \textbf{Name} & \textbf{Meaning} & \textbf{Driver} \\
\midrule
BR & Biological Response & True pharmacological effect & Treatment, biomarker \\
ER & Expectancy Response & Placebo/expectancy effect & Blinding status \\
TR & Time-variant Response & Natural disease progression & Time, individual \\
\bottomrule
\end{tabular}
\end{table}

\section{Component Mean Formulas}

\subsection{Biological Response (BR)}

\begin{equation}
\text{BR}_{mean,t} = \begin{cases}
w_t \cdot r_{eff} & \text{if on treatment} \\
w_{accum} \cdot \gamma \cdot f(t_{off}) & \text{if off treatment (carryover)}
\end{cases}
\end{equation}

where:
\begin{itemize}
    \item $w_t$ = weeks on drug up to timepoint $t$
    \item $r_{eff} = r_0 (1 + \beta \cdot \text{BiomarkerCentered})$ = effective BR rate
    \item $\gamma$ = carryover proportion
    \item $f(t_{off})$ = decay function (exponential, power law, etc.)
\end{itemize}

\subsection{Expectancy Response (ER)}

\begin{equation}
\text{ER}_{mean,t} = e_t \cdot r_{exp}
\end{equation}

where:
\begin{itemize}
    \item $e_t \in \{0, 0.5, 1\}$ = expectancy factor (design-dependent)
    \item $r_{exp}$ = expectancy response rate
\end{itemize}

\subsection{Time-variant Response (TR)}

\begin{equation}
\text{TR}_{mean,t} = f_{time}(t)
\end{equation}

representing natural disease trajectory independent of treatment.

\section{Design-Specific Expectancy Patterns}

\begin{table}[h]
\centering
\caption{Expectancy Allocation by Design}
\begin{tabular}{@{}lll@{}}
\toprule
\textbf{Design} & \textbf{Weeks with $e=1.0$} & \textbf{Rationale} \\
\midrule
OL+BDC & Weeks $\leq 16$ & Open-label phase \\
Hybrid & Weeks 4, 8 only & Unblinded baseline only \\
Crossover & Varies by period & Blinding maintained \\
\bottomrule
\end{tabular}
\end{table}

%==============================================================================
\part{Biomarker$\times$Treatment Interaction}
%==============================================================================

\section{Two-Level Mechanism}

The biomarker-treatment interaction operates at two levels:

\subsection{Level 1: Covariance Structure ($c_{bm}$)}

The correlation between biomarker and BR components:
\begin{equation}
\text{Corr}(\text{biomarker}, \text{BR}_t) = c_{bm}
\end{equation}

This creates covariance-level association: participants with high biomarker
values tend to have correlated BR values.

\subsection{Level 2: Mean Modulation (biomarker\_moderation)}

The biomarker modifies treatment effect magnitude:
\begin{equation}
r_{eff} = r_0 (1 + \beta \cdot \text{BiomarkerCentered})
\end{equation}

where $\beta$ = \texttt{biomarker\_moderation} parameter.

\section{Mathematical Relationship}

\begin{table}[h]
\centering
\caption{Comparison of Biomarker Effect Mechanisms}
\begin{tabular}{@{}lll@{}}
\toprule
\textbf{Aspect} & \textbf{$c_{bm}$} & \textbf{biomarker\_moderation} \\
\midrule
Level & Covariance & Mean \\
Effect & Correlation in draws & Scaling of treatment effect \\
Parameter range & 0--0.6 (PD constraint) & 0--0.65 (power analysis) \\
Impact & Noise structure & Signal magnitude \\
\bottomrule
\end{tabular}
\end{table}

\section{Detectability}

The biomarker$\times$treatment interaction is tested via:
\begin{equation}
H_0: \beta_{biomarker \times treatment} = 0
\end{equation}

in the mixed model:
\begin{lstlisting}
lmer(response ~ biomarker * treatment + week + (1|participant_id))
\end{lstlisting}

Power depends on both $c_{bm}$ and biomarker\_moderation values.

%==============================================================================
\part{Carryover Effect Modeling}
%==============================================================================

\section{The Fundamental Question}

Should carryover effects modify correlation parameters in addition to means?

\begin{proposition}
Carryover affects means only, not correlations.
\end{proposition}

\section{Mathematical Argument}

\textbf{Approach 1 (Incorrect)}: Carryover modifies correlations
\begin{equation}
\rho_{BR}(t_{on}, t_{off}) = \rho_0 \cdot g(\text{carryover})
\end{equation}

This would mean participants' BR correlations change based on their treatment
history---biologically implausible for most drugs.

\textbf{Approach 2 (Correct)}: Carryover affects means only
\begin{equation}
\mu_{BR}(t_{off}) = \mu_{BR}(t_{on}) \cdot h(t_{since\_discontinuation})
\end{equation}

The residual drug effect enters through the expected value, not the variance
structure.

\section{Double-Counting Problem}

If carryover modified both means and correlations:
\begin{enumerate}
    \item Mean effect: $\text{BR}_{mean}$ higher during off-period
    \item Correlation effect: BR draws more correlated during off-period
\end{enumerate}

This would count the same biological phenomenon twice, inflating carryover
impact and biasing estimates.

\section{Five Carryover Models}

\begin{table}[h]
\centering
\caption{Carryover Model Comparison}
\begin{tabular}{@{}llll@{}}
\toprule
\textbf{Model} & \textbf{Formula} & \textbf{Parameters} & \textbf{Best For} \\
\midrule
Fixed Proportion & $\gamma$ (constant) & 1 & Simple drugs \\
Exponential & $\exp(-\lambda t)$ & 2 & Psychiatric drugs \\
Power Law & $(1+\alpha t)^{-\beta}$ & 3 & Tissue accumulation \\
Two-Compartment & $Ae^{-\alpha t} + Be^{-\beta t}$ & 4 & Biphasic kinetics \\
Delayed Elimination & Lag + exponential & 3 & Prodrugs \\
\bottomrule
\end{tabular}
\end{table}

\section{Half-Life Relationship}

For exponential decay:
\begin{equation}
t_{1/2} = \frac{\ln(2)}{\lambda} \approx \frac{0.693}{\lambda}
\end{equation}

\begin{table}[h]
\centering
\caption{Half-Life Examples}
\begin{tabular}{@{}lll@{}}
\toprule
\textbf{$\lambda$ (week$^{-1}$)} & \textbf{$t_{1/2}$ (weeks)} & \textbf{Application} \\
\midrule
0.05 & 13.9 & Slow-elimination drugs \\
0.10 & 6.9 & Typical psychiatric medications \\
0.15 & 4.6 & Faster clearance \\
\bottomrule
\end{tabular}
\end{table}

%==============================================================================
\part{Mixed-Effects Model Specification}
%==============================================================================

\section{Core Model Formula}

\begin{align}
\text{Response}_{it} &= \beta_0 + \beta_1 \cdot \text{Biomarker}_i +
\beta_2 \cdot \text{Treatment}_{it} \notag \\
&\quad + \beta_3 \cdot \text{Biomarker}_i \times \text{Treatment}_{it} +
\beta_4 \cdot \text{Week}_{it} + u_i + \epsilon_{it}
\end{align}

where:
\begin{itemize}
    \item $\beta_3$ = biomarker$\times$treatment interaction (primary endpoint)
    \item $u_i \sim N(0, \sigma_u^2)$ = random intercept
    \item $\epsilon_{it} \sim N(0, \sigma^2)$ = residual error
\end{itemize}

\section{Model Variants}

\textbf{With Carryover Term}:
\begin{lstlisting}
lmer(response ~ biomarker * treatment + week + carryover_effect +
     (1|participant_id))
\end{lstlisting}

\textbf{Without Carryover (Hendrickson Original)}:
\begin{lstlisting}
lmer(response ~ biomarker * treatment + week + (1|participant_id))
\end{lstlisting}

\section{Random Effects Options}

\begin{table}[h]
\centering
\caption{Random Effects Structures}
\begin{tabular}{@{}lll@{}}
\toprule
\textbf{Structure} & \textbf{Formula} & \textbf{Interpretation} \\
\midrule
Random intercept & \texttt{(1|participant\_id)} & Individual baseline differences \\
Random slope (drug) & \texttt{(1+treatment|participant\_id)} & Individual treatment response \\
Random slope (time) & \texttt{(1+week|participant\_id)} & Individual trajectories \\
\bottomrule
\end{tabular}
\end{table}

%==============================================================================
\part{Design Comparison Framework}
%==============================================================================

\section{Five Trial Designs}

\begin{table}[h]
\centering
\caption{Trial Design Comparison}
\begin{tabular}{@{}lllll@{}}
\toprule
\textbf{Design} & \textbf{Points} & \textbf{Schedule} & \textbf{Carryover} \\
\midrule
Hybrid & 8 & Clustered & Moderate-High \\
OL+BDC & 8 & Clustered & High \\
Open-Label & 4 & Evenly-spaced & None \\
Crossover & 4 & Evenly-spaced & Critical \\
Parallel & 4 & Evenly-spaced & None \\
\bottomrule
\end{tabular}
\end{table}

\section{Measurement Schedules}

\subsection{Clustered (Hybrid)}

Weeks: $\{4, 8, 9, 10, 11, 12, 16, 20\}$
\begin{itemize}
    \item Dense cluster at transition (weeks 9--12)
    \item Captures carryover decay dynamics
    \item AR(1) works well ($\kappa \approx 56$--$61$)
\end{itemize}

\subsection{Clustered (OL+BDC)}

Weeks: $\{4, 8, 12, 16, 17, 18, 19, 20\}$
\begin{itemize}
    \item Dense cluster at discontinuation (weeks 16--20)
    \item Extended washout observation
    \item AR(1) works well ($\kappa \approx 60$)
\end{itemize}

\subsection{Evenly-Spaced}

Weeks: $\{2, 8, 14, 20\}$
\begin{itemize}
    \item Uniform 6-week gaps
    \item Requires alternative correlation structures
    \item Exponential or Power Law recommended ($\kappa \approx 13$--$18$)
\end{itemize}

\section{Expectancy Vulnerability Analysis}

\textbf{OL+BDC}: MORE vulnerable to expectancy confounding
\begin{itemize}
    \item High expectancy ($e=1.0$) for 80\% of trial (weeks 0--16)
    \item ER and BR effects coupled at randomization point
    \item Sudden drop in both components at week 16
\end{itemize}

\textbf{Hybrid}: LESS vulnerable
\begin{itemize}
    \item High expectancy only for 20\% of trial (weeks 4--8)
    \item Most measurements during blinded phase ($e=0.5$)
    \item Cleaner separation of BR and ER effects
\end{itemize}

%==============================================================================
\part{Statistical Power Analysis}
%==============================================================================

\section{Power Definition}

\begin{equation}
\text{Power} = P(\text{Reject } H_0 | H_1 \text{ true})
\end{equation}

Estimated via Monte Carlo simulation:
\begin{equation}
\hat{\text{Power}} = \frac{\#(\text{iterations with } p < 0.05)}{N_{iterations}}
\end{equation}

\section{Power Determinants (Ranked by Impact)}

\begin{table}[h]
\centering
\caption{Parameter Impact Ranking}
\begin{tabular}{@{}lll@{}}
\toprule
\textbf{Rank} & \textbf{Parameter} & \textbf{Impact Mechanism} \\
\midrule
1 & Design & Measurement schedule, carryover modeling \\
2 & biomarker\_moderation & Signal magnitude \\
3 & carryover\_t1half & Signal clarity during off-periods \\
4 & $c_{bm}$ & Biomarker-response correlation \\
5 & n\_participants & Statistical precision \\
6 & Correlation structure & Numerical stability \\
\bottomrule
\end{tabular}
\end{table}

\section{Carryover Cost Analysis}

\begin{table}[h]
\centering
\caption{Power Loss from Ignoring Carryover}
\begin{tabular}{@{}llll@{}}
\toprule
\textbf{Carryover} & \textbf{Without Term} & \textbf{With Term} & \textbf{Type I Error} \\
\midrule
$\gamma = 0$ & $\sim$0.50 & $\sim$0.50 & 0.05 \\
$\gamma = 0.5$ & $\sim$0.35 & $\sim$0.45 & 0.07--0.10 (without) \\
$\gamma = 1.0$ & $\sim$0.25 & $\sim$0.35 & 0.10--0.12 (without) \\
\bottomrule
\end{tabular}
\end{table}

Ignoring carryover leads to:
\begin{itemize}
    \item 10--20\% power reduction
    \item Inflated Type I error
    \item Biased treatment effect estimates
\end{itemize}

%==============================================================================
\part{Alternative Summary Statistics for Interaction Testing}
%==============================================================================

\section{Motivation}

The primary analysis in this simulation framework uses mixed-effects models to
test the biomarker$\times$treatment interaction:
\begin{lstlisting}
lmer(response ~ biomarker * treatment + week + (1|participant_id))
\end{lstlisting}

While this approach is flexible and accounts for repeated measures, the test
statistic (Wald $z$ or $t$) lacks a simple closed-form power expression. This
section develops \textbf{alternative summary statistics} that:

\begin{enumerate}
    \item Permit closed-form power calculations
    \item Provide intuition for optimal trial design
    \item May sacrifice some efficiency for analytical tractability
\end{enumerate}

\section{The Correlation-Based Test Statistic}

\subsection{Core Intuition}

The biomarker$\times$treatment interaction implies that participants with
higher biomarker values respond more strongly to treatment. This manifests as
a correlation between:

\begin{itemize}
    \item \textbf{Biomarker level}: $X_i$ (baseline covariate)
    \item \textbf{Treatment response differential}: $\Delta Y_i = \bar{Y}_{i,on} - \bar{Y}_{i,off}$
\end{itemize}

Under the null hypothesis of no interaction, this correlation should be zero.
Under the alternative, it should be positive (assuming higher biomarker
implies stronger treatment effect).

\subsection{Definition}

Define the test statistic:
\begin{equation}
\boxed{r = \text{Corr}(X_i, \Delta Y_i)}
\end{equation}

where:
\begin{itemize}
    \item $X_i$ = centered biomarker value for participant $i$
    \item $\bar{Y}_{i,on}$ = mean response when participant $i$ is on treatment
    \item $\bar{Y}_{i,off}$ = mean response when participant $i$ is off treatment
    \item $\Delta Y_i = \bar{Y}_{i,on} - \bar{Y}_{i,off}$ = within-participant treatment effect
\end{itemize}

\subsection{Computation}

For each participant $i = 1, \ldots, n$:
\begin{align}
\bar{Y}_{i,on} &= \frac{1}{n_{on,i}} \sum_{t: T_{it}=1} Y_{it} \\
\bar{Y}_{i,off} &= \frac{1}{n_{off,i}} \sum_{t: T_{it}=0} Y_{it} \\
\Delta Y_i &= \bar{Y}_{i,on} - \bar{Y}_{i,off}
\end{align}

Then compute the sample correlation:
\begin{equation}
r = \frac{\sum_{i=1}^{n}(X_i - \bar{X})(\Delta Y_i - \overline{\Delta Y})}
{\sqrt{\sum_{i=1}^{n}(X_i - \bar{X})^2 \sum_{i=1}^{n}(\Delta Y_i - \overline{\Delta Y})^2}}
\end{equation}

\section{Distribution Under the Null}

\subsection{Fisher's $z$-Transformation}

Under $H_0: \rho = 0$ (no biomarker$\times$treatment interaction):
\begin{equation}
z = \frac{1}{2} \ln\left(\frac{1+r}{1-r}\right) \approx \mathcal{N}\left(0, \frac{1}{n-3}\right)
\end{equation}

This transformation stabilizes the variance, making it approximately
independent of the true correlation.

\subsection{Test Procedure}

Reject $H_0$ at level $\alpha$ if:
\begin{equation}
|z| > z_{1-\alpha/2} / \sqrt{n-3}
\end{equation}

or equivalently, if:
\begin{equation}
|r| > \frac{z_{1-\alpha/2}}{\sqrt{n-3 + z_{1-\alpha/2}^2}}
\end{equation}

For $\alpha = 0.05$ and $n = 70$:
\begin{equation}
|r| > \frac{1.96}{\sqrt{67 + 3.84}} \approx 0.233
\end{equation}

\section{Distribution Under the Alternative}

\subsection{True Correlation}

Under $H_1$, the true correlation $\rho$ between biomarker and treatment
differential depends on:

\begin{equation}
\rho = \frac{\text{Cov}(X_i, \Delta Y_i)}{\sqrt{\text{Var}(X_i) \cdot \text{Var}(\Delta Y_i)}}
\end{equation}

\subsection{Deriving $\text{Cov}(X_i, \Delta Y_i)$}

Under the linear interaction model:
\begin{equation}
E[\Delta Y_i | X_i] = \tau + \beta \cdot X_i
\end{equation}

where $\beta$ is the biomarker$\times$treatment interaction coefficient.
Therefore:
\begin{equation}
\text{Cov}(X_i, \Delta Y_i) = \beta \cdot \text{Var}(X_i)
\end{equation}

\subsection{Deriving $\text{Var}(\Delta Y_i)$}

The variance of the treatment differential has two sources:
\begin{align}
\text{Var}(\Delta Y_i) &= \beta^2 \cdot \text{Var}(X_i) + \text{Var}(\epsilon_{\Delta,i}) \\
&= \beta^2 \sigma_X^2 + \sigma_{\Delta}^2
\end{align}

where $\sigma_{\Delta}^2$ is the residual variance of the treatment effect
(reflecting measurement error, within-person variability, etc.).

\subsection{Closed-Form Expression for $\rho$}

Combining these:
\begin{equation}
\boxed{\rho = \frac{\beta \sigma_X}{\sqrt{\beta^2 \sigma_X^2 + \sigma_{\Delta}^2}}
= \frac{1}{\sqrt{1 + \frac{\sigma_{\Delta}^2}{\beta^2 \sigma_X^2}}}}
\end{equation}

This has a signal-to-noise interpretation:
\begin{equation}
\rho = \frac{1}{\sqrt{1 + 1/\text{SNR}^2}} \quad \text{where} \quad
\text{SNR} = \frac{\beta \sigma_X}{\sigma_{\Delta}}
\end{equation}

\section{Closed-Form Power Expression}

\subsection{Power via Fisher's $z$}

Under the alternative with true correlation $\rho$:
\begin{equation}
z \sim \mathcal{N}\left(\zeta, \frac{1}{n-3}\right)
\end{equation}

where $\zeta = \frac{1}{2}\ln\left(\frac{1+\rho}{1-\rho}\right)$ is the
population Fisher's $z$.

\subsection{Power Formula}

\begin{theorem}[Power for Correlation Test]
The power to detect a biomarker$\times$treatment interaction with true
correlation $\rho$ using $n$ participants at significance level $\alpha$ is:
\begin{equation}
\boxed{1 - \beta = \Phi\left(\zeta\sqrt{n-3} - z_{1-\alpha/2}\right)
+ \Phi\left(-\zeta\sqrt{n-3} - z_{1-\alpha/2}\right)}
\end{equation}

For one-sided tests (expected positive correlation):
\begin{equation}
1 - \beta = \Phi\left(\zeta\sqrt{n-3} - z_{1-\alpha}\right)
\end{equation}
\end{theorem}

\subsection{Sample Size Formula}

Inverting the power formula for required sample size:
\begin{equation}
\boxed{n = 3 + \left(\frac{z_{1-\alpha/2} + z_{1-\beta}}{\zeta}\right)^2}
\end{equation}

\begin{table}[h]
\centering
\caption{Required Sample Size for 80\% Power ($\alpha = 0.05$, two-sided)}
\begin{tabular}{@{}lll@{}}
\toprule
\textbf{True $\rho$} & \textbf{Fisher's $\zeta$} & \textbf{Required $n$} \\
\midrule
0.20 & 0.203 & 197 \\
0.30 & 0.310 & 88 \\
0.40 & 0.424 & 49 \\
0.50 & 0.549 & 32 \\
0.60 & 0.693 & 23 \\
\bottomrule
\end{tabular}
\end{table}

\section{Relationship to Interaction Coefficient}

\subsection{Converting Between $\rho$ and $\beta$}

The regression coefficient $\beta$ (biomarker\_moderation) and correlation
$\rho$ are related by:
\begin{equation}
\beta = \rho \cdot \frac{\sigma_{\Delta}}{\sigma_X}
\end{equation}

or equivalently:
\begin{equation}
\rho = \beta \cdot \frac{\sigma_X}{\sigma_{\Delta}}
\end{equation}

\subsection{Interpretation}

\begin{itemize}
    \item $\beta$ measures the \textbf{unstandardized effect}: how many units
    the treatment differential changes per unit increase in biomarker
    \item $\rho$ measures the \textbf{standardized association}: the
    correlation between biomarker and treatment response
\end{itemize}

For power calculations, $\rho$ is more useful because its distribution
is known (via Fisher's $z$).

\section{Design Implications}

\subsection{What Determines $\rho$?}

From the SNR formula, larger $\rho$ (easier detection) occurs when:

\begin{enumerate}
    \item \textbf{Larger $\beta$}: Stronger true interaction effect
    \item \textbf{Larger $\sigma_X$}: More biomarker heterogeneity in sample
    \item \textbf{Smaller $\sigma_{\Delta}$}: Less noise in treatment differential
\end{enumerate}

\subsection{Reducing $\sigma_{\Delta}$}

The variance of the treatment differential can be decomposed:
\begin{equation}
\sigma_{\Delta}^2 = \frac{\sigma_{on}^2}{n_{on}} + \frac{\sigma_{off}^2}{n_{off}}
+ 2\sigma_{\text{within}}^2
\end{equation}

where:
\begin{itemize}
    \item $\sigma_{on}^2, \sigma_{off}^2$ = measurement error variance
    \item $n_{on}, n_{off}$ = number of measurements in each condition
    \item $\sigma_{\text{within}}^2$ = within-person variability
\end{itemize}

\textbf{Design implication}: More measurements per participant reduce
$\sigma_{\Delta}^2$, increasing $\rho$ and power.

\subsection{Balance Between Conditions}

For fixed total measurements $n_{on} + n_{off} = m$, $\sigma_{\Delta}^2$ is
minimized when $n_{on} = n_{off} = m/2$ (balanced allocation).

\textbf{Design implication}: Balanced designs (equal on/off measurements)
maximize power for detecting biomarker$\times$treatment interactions.

\subsection{Recruitment Strategy}

Power increases with $\sigma_X$ (biomarker variance). This suggests:

\begin{enumerate}
    \item Recruit from heterogeneous populations
    \item Avoid restricting biomarker range
    \item Consider stratified sampling to ensure biomarker spread
\end{enumerate}

\section{Comparison to Mixed-Effects Approach}

\begin{table}[h]
\centering
\caption{Comparison of Interaction Test Approaches}
\begin{tabular}{@{}lll@{}}
\toprule
\textbf{Aspect} & \textbf{Correlation Test} & \textbf{Mixed Model} \\
\midrule
Test statistic & $r$ (or Fisher's $z$) & Wald $z$ or $t$ \\
Power formula & Closed-form (Fisher) & Monte Carlo only \\
Repeated measures & Collapsed to $\Delta Y$ & Modeled explicitly \\
Covariates & Not directly included & Easily added \\
Efficiency & Lower (loses information) & Higher (uses all data) \\
Robustness & Less model-dependent & Assumes linearity, normality \\
\bottomrule
\end{tabular}
\end{table}

\subsection{When to Use Each}

\textbf{Correlation test}:
\begin{itemize}
    \item Power/sample size planning (closed-form calculations)
    \item Quick exploratory analysis
    \item When model assumptions are uncertain
\end{itemize}

\textbf{Mixed model}:
\begin{itemize}
    \item Primary confirmatory analysis
    \item When covariates must be controlled
    \item When repeated measures structure is important
\end{itemize}

\section{Extension: Slope-Based Statistics}

\subsection{Regression of $\Delta Y$ on $X$}

An alternative to correlation is the regression slope:
\begin{equation}
\Delta Y_i = \alpha + \beta X_i + \epsilon_i
\end{equation}

The OLS estimate:
\begin{equation}
\hat{\beta} = \frac{\sum_{i}(X_i - \bar{X})(\Delta Y_i - \overline{\Delta Y})}
{\sum_{i}(X_i - \bar{X})^2}
\end{equation}

Under standard assumptions:
\begin{equation}
\hat{\beta} \sim \mathcal{N}\left(\beta, \frac{\sigma_{\Delta}^2}{\sum_{i}(X_i - \bar{X})^2}\right)
\end{equation}

\subsection{Power for Slope Test}

\begin{equation}
1 - \beta = \Phi\left(\frac{|\beta|\sqrt{\sum(X_i - \bar{X})^2}}{\sigma_{\Delta}}
- z_{1-\alpha/2}\right)
\end{equation}

This is equivalent to the correlation test but expressed in terms of the
regression coefficient rather than the correlation.

\section{Summary of Key Formulas}

\subsection{Test Statistic}
\begin{equation}
r = \text{Corr}(X_i, \Delta Y_i) = \text{Corr}(\text{Biomarker}, \bar{Y}_{on} - \bar{Y}_{off})
\end{equation}

\subsection{Signal-to-Noise Ratio}
\begin{equation}
\rho = \frac{1}{\sqrt{1 + \sigma_{\Delta}^2 / (\beta^2 \sigma_X^2)}}
\end{equation}

\subsection{Power}
\begin{equation}
1 - \beta = \Phi\left(\zeta\sqrt{n-3} - z_{1-\alpha/2}\right)
\end{equation}

\subsection{Sample Size}
\begin{equation}
n = 3 + \left(\frac{z_{1-\alpha/2} + z_{1-\beta}}{\zeta}\right)^2
\end{equation}

\subsection{Design Principle}
\begin{quote}
\textit{To maximize power for detecting biomarker$\times$treatment interactions:
maximize biomarker heterogeneity ($\sigma_X$), increase measurements per
condition ($n_{on}, n_{off}$), and maintain balanced allocation.}
\end{quote}

\section{Literature Review: Analytical Power Alternatives}

This section surveys analytical approaches to power calculation in settings
comparable to Hendrickson et al.\ (2020), where mixed-effects models test
biomarker$\times$treatment interactions in N-of-1 or crossover designs.

\subsection{Why Analytical Power Is Difficult}

The mixed-effects model
\begin{lstlisting}
lmer(response ~ biomarker * treatment + week + (1|participant_id))
\end{lstlisting}
involves a Wald test statistic whose exact distribution depends on:
\begin{enumerate}
    \item The design matrix structure $X$
    \item The variance-covariance matrix $V$
    \item The variance components ($\sigma_u^2$, $\sigma^2$)
\end{enumerate}

Computing $\text{Var}(\hat{\beta}) = \sigma^2(X'V^{-1}X)^{-1}$ requires inverting
a large block-diagonal matrix whose structure varies by design. For Hendrickson's
clustered measurement schedules with AR(1) correlation, this inversion lacks a
simple closed form.

\subsection{Summary Measures Analysis}

The most widely used analytical alternative collapses repeated measures to
participant-level summaries.

\textbf{Key references:}
\begin{itemize}
    \item Matthews, Altman, Campbell, Royston (1990) -- ``Analysis of serial
    measurements in medical research'' (BMJ)
    \item Frison \& Pocock (1992) -- ``Repeated measures in clinical trials:
    analysis using mean summary statistics''
\end{itemize}

\textbf{Approach:} Compute a summary statistic per participant (mean, AUC,
slope, change score), then apply standard between-subject tests. The correlation
$r = \text{Corr}(X_i, \Delta Y_i)$ developed in this document is a summary
measure approach.

\textbf{Trade-off:} Loses efficiency (ignores within-person correlation structure)
but gains analytical tractability. Efficiency loss is moderate when intraclass
correlation (ICC) is high (0.5--0.8), which is typical in N-of-1 settings.

\subsection{The ``Rule of $4\times$'' for Interactions}

\textbf{Key reference:} Brookes et al.\ (2004) -- ``Subgroup analyses in
randomised trials: risks of subgroup-specific analyses'' (Statistics in Medicine)

\textbf{Key finding:} Detecting a treatment$\times$covariate interaction requires
approximately \textbf{4 times the sample size} compared to detecting the main
treatment effect of the same magnitude.

\textbf{Mathematical basis:}
\begin{equation}
\text{Var}(\hat{\beta}_{interaction}) \approx 4 \cdot \text{Var}(\hat{\beta}_{treatment})
\end{equation}

This arises because the interaction term has reduced ``effective sample size''---only
the covariance between biomarker and treatment contributes to the signal.

\textbf{Implication:} If detecting a main treatment effect requires $n = 20$,
detecting the same-sized interaction requires $n \approx 80$.

\subsection{Asymptotic Mixed Model Formulas}

\textbf{Key references:}
\begin{itemize}
    \item Shieh (2003) -- ``On power and sample size calculations for likelihood
    ratio tests in generalized linear models''
    \item Stroup (2002) -- ``Power analysis based on spatial effects mixed models''
    \item Demidenko (2004) -- \textit{Mixed Models: Theory and Applications} (Chapter 10)
\end{itemize}

\textbf{Approach:} Use the noncentrality parameter of the Wald or likelihood
ratio test:
\begin{equation}
\text{ncp} = \frac{\beta^2}{\text{Var}(\hat{\beta})} =
\frac{\beta^2}{\sigma^2 (X'V^{-1}X)^{-1}_{jj}}
\end{equation}

Then:
\begin{equation}
\text{Power} = \Phi\left(\frac{|\beta|}{\text{SE}(\hat{\beta})} - z_{1-\alpha/2}\right)
\end{equation}

\textbf{Challenge:} Computing $(X'V^{-1}X)^{-1}$ requires specifying the design
matrix, variance components, and correlation structure. For balanced designs
with compound symmetry, closed forms exist. For Hendrickson's clustered designs
with AR(1), the algebra becomes intractable.

\subsection{GEE-Based Power}

\textbf{Key references:}
\begin{itemize}
    \item Liu \& Liang (1997) -- ``Sample size calculations for studies with
    correlated observations''
    \item Rochon (1998) -- ``Application of GEE procedures for sample size
    calculations in repeated measures experiments''
\end{itemize}

\textbf{Formula for treatment effect:}
\begin{equation}
n = \frac{2(z_{\alpha/2} + z_{\beta})^2 \sigma^2}{\delta^2} \times \text{Design Effect}
\end{equation}

where Design Effect $= \frac{1 + (m-1)\rho}{m}$ for compound symmetry with $m$
measurements and correlation $\rho$.

\textbf{Limitation:} Primarily developed for main effects, not
treatment$\times$covariate interactions.

\subsection{Crossover Trial Formulas}

\textbf{Key references:}
\begin{itemize}
    \item Jones \& Kenward (2014) -- \textit{Design and Analysis of Cross-Over
    Trials} (3rd edition)
    \item Senn (2002) -- \textit{Cross-over Trials in Clinical Research} (2nd edition)
\end{itemize}

\textbf{For standard $2 \times 2$ crossover (treatment effect):}
\begin{equation}
n = \frac{2(z_{\alpha/2} + z_{\beta})^2 \sigma_d^2}{\delta^2}
\end{equation}

where $\sigma_d^2$ is the within-subject variance of differences.

\textbf{For treatment$\times$period interaction} (analogous to
treatment$\times$biomarker):
\begin{equation}
n_{interaction} \approx 4 \times n_{treatment}
\end{equation}

\textbf{Limitation:} These formulas assume balanced designs. Hendrickson's Hybrid
and OL+BDC designs are unbalanced (different numbers of on/off measurements).

\subsection{Design Effect Approach}

\textbf{Key reference:} Diggle, Heagerty, Liang, Zeger (2002) --
\textit{Analysis of Longitudinal Data} (Chapter 13)

\textbf{Approach:} Express power in terms of a ``design effect'' that inflates variance:
\begin{equation}
\text{Var}(\hat{\beta}) = \text{Var}_{simple} \times DE
\end{equation}

For AR(1) correlation with $m$ equally-spaced timepoints:
\begin{equation}
DE = \frac{1}{m} \left[1 + 2\sum_{k=1}^{m-1}\left(1 - \frac{k}{m}\right)\rho^k\right]
\end{equation}

\textbf{For Hendrickson's clustered designs:} The unequal spacing
(weeks 4, 8, 9, 10, 11, 12, 16, 20) requires computing:
\begin{equation}
DE = \frac{1}{m^2} \sum_{i,j} \rho^{|t_i - t_j|}
\end{equation}

This can be computed numerically but does not yield a simple closed form.

\subsection{Bayesian Assurance}

\textbf{Key reference:} O'Hagan \& Stevens (2001) -- ``Bayesian assessment of
sample size for clinical trials of cost-effectiveness''

\textbf{Approach:} Instead of frequentist power (probability of significance
given fixed true effect), compute ``assurance'' (probability of significance
averaging over prior uncertainty about true effect):
\begin{equation}
\text{Assurance} = \int \text{Power}(\theta) \cdot \pi(\theta) \, d\theta
\end{equation}

\textbf{Advantage:} More realistic when effect size is uncertain.

\textbf{Disadvantage:} Requires specifying a prior; still often computed by
simulation.

\subsection{Comparison of Analytical Approaches}

\begin{table}[h]
\centering
\caption{Comparison of Power Calculation Methods}
\begin{tabular}{@{}lcccc@{}}
\toprule
\textbf{Approach} & \textbf{Analytical?} & \textbf{Interaction?} &
\textbf{Unbalanced?} & \textbf{Accuracy} \\
\midrule
Summary measures & Yes & Yes & Partial & Moderate \\
Rule of $4\times$ & Yes & Yes & No & Approximate \\
Asymptotic mixed & Semi & Yes & Numerical & Good \\
GEE-based & Yes & Limited & Partial & Moderate \\
Crossover formulas & Yes & Limited & No & Good (balanced) \\
Design effect & Semi & Limited & Yes & Moderate \\
Simulation & No & Yes & Yes & Excellent \\
\bottomrule
\end{tabular}
\end{table}

\subsection{The Efficiency-Tractability Trade-off}

Three approaches represent different points on the efficiency-tractability
spectrum:

\begin{enumerate}
    \item \textbf{Full mixed model} (Hendrickson approach): Uses all data optimally
    but requires simulation for power.

    \item \textbf{Summary correlation} (this document): Loses $\sim$15--30\%
    efficiency but has closed-form power via Fisher's $z$.

    \item \textbf{Rule of $4\times$}: Very rough approximation---if main effect
    needs $n$, interaction needs $4n$.
\end{enumerate}

The summary approach is well-established in the literature (Matthews et al., 1990).
It is particularly appropriate when ICC is high (typical in N-of-1: 0.5--0.8),
where collapsing measurements loses minimal information.

\subsection{Recommended Analytical Approximation}

For Hendrickson-style power calculations, a practical analytical formula combines:

\textbf{Step 1:} Estimate effective correlation from design parameters
\begin{equation}
\rho_{eff} = \frac{\beta \cdot \sigma_X}{\sqrt{\beta^2 \sigma_X^2 + \sigma_{\Delta}^2}}
\end{equation}

where $\sigma_{\Delta}^2 \approx \frac{\sigma^2}{n_{on}} + \frac{\sigma^2}{n_{off}}$
(measurement noise reduced by averaging).

\textbf{Step 2:} Apply Fisher's $z$ power formula
\begin{equation}
n = 3 + \left(\frac{z_{1-\alpha/2} + z_{1-\beta}}{\zeta}\right)^2
\end{equation}

\textbf{Step 3:} Apply design-specific inflation factor

\begin{table}[h]
\centering
\caption{Design-Specific Inflation Factors}
\begin{tabular}{@{}lll@{}}
\toprule
\textbf{Design} & \textbf{Inflation} & \textbf{Rationale} \\
\midrule
Hybrid & $\times 1.1$ & Near-balanced on/off measurements \\
OL+BDC & $\times 1.2$ & Unbalanced (more on than off) \\
Crossover & $\times 1.1$ & Balanced but carryover concerns \\
Parallel & $\times 1.5$ & Between-subject comparison loses power \\
\bottomrule
\end{tabular}
\end{table}

This hybrid approach provides analytical tractability while acknowledging
design-specific features through calibrated inflation factors.

\subsection{Key References Summary}

\begin{enumerate}
    \item Matthews et al.\ (1990) -- Summary measures analysis (BMJ)
    \item Frison \& Pocock (1992) -- Mean summary statistics
    \item Brookes et al.\ (2004) -- Rule of $4\times$ for interactions
    \item Shieh (2003) -- Asymptotic mixed model power
    \item Liu \& Liang (1997) -- GEE-based power
    \item Jones \& Kenward (2014) -- Crossover trial design
    \item Senn (2002) -- Cross-over trials in clinical research
    \item Diggle et al.\ (2002) -- Analysis of longitudinal data
    \item Demidenko (2004) -- Mixed models theory
    \item O'Hagan \& Stevens (2001) -- Bayesian assurance
\end{enumerate}

%==============================================================================
\part{Key Mathematical Insights}
%==============================================================================

\section{The Eigenvalue-Conditioning Nexus}

The condition number problem is fundamentally about eigenvalue balance:
\begin{equation}
\kappa = \frac{\lambda_{max}}{\lambda_{min}}
\end{equation}

\textbf{Good conditioning} requires eigenvalues to be similar in magnitude,
which requires correlation values to span the correlation space evenly rather
than clustering at extreme values.

\section{The Clustering Insight}

Clustered measurement schedules succeed because they create
\textbf{intermediate correlations}:
\begin{itemize}
    \item AR(1) with 1-week gaps: $\rho = 0.75$
    \item AR(1) with 4-week gaps: $\rho = 0.32$
    \item AR(1) with 8-week gaps: $\rho = 0.10$
\end{itemize}

This natural variation from clustering balances eigenvalues.

Evenly-spaced designs with uniform large gaps (6 weeks) create extreme
correlation ratios (32:1), leading to extreme eigenvalue ratios.

\section{The Carryover-Correlation Separation}

The principle that carryover affects means only (not correlations) is
essential for:
\begin{enumerate}
    \item \textbf{Avoiding double-counting} of carryover effects
    \item \textbf{Maintaining interpretability} of correlation parameters
    \item \textbf{Preserving positive definiteness}
\end{enumerate}

\section{The Hierarchy Constraint}

The correlation hierarchy ($c_{cfct} < c_{cf1t} < c_{autocorr}$) encodes
the temporal coherence principle:

\begin{quote}
\textit{Measurements at different times cannot be more correlated than
measurements at the same time.}
\end{quote}

Violating this creates mathematically impossible constraint systems.

%==============================================================================
\appendix
\part*{Appendices}
\addcontentsline{toc}{part}{Appendices}
%==============================================================================

\section{Key Formulas Reference}

\subsection{Covariance Matrix}
\begin{equation}
\Sigma = D \cdot R \cdot D
\end{equation}

\subsection{Conditional Mean}
\begin{equation}
\mu_{1|2} = \mu_1 + \Sigma_{12} \Sigma_{22}^{-1} (x_2 - \mu_2)
\end{equation}

\subsection{Conditional Covariance}
\begin{equation}
\Sigma_{cond} = \Sigma_{11} - \Sigma_{12} \Sigma_{22}^{-1} \Sigma_{12}'
\end{equation}

\subsection{Condition Number}
\begin{equation}
\kappa = \frac{\lambda_{max}}{\lambda_{min}}
\end{equation}

\subsection{AR(1) Correlation}
\begin{equation}
\rho_{ij} = \rho^{|t_i - t_j|}
\end{equation}

\subsection{Exponential Correlation}
\begin{equation}
\rho_{ij} = \exp(-\lambda |t_i - t_j|)
\end{equation}

\subsection{Carryover Exponential Decay}
\begin{equation}
\text{Carryover}_t = \text{Effect}_0 \cdot \exp\left(-\frac{\ln(2)}{t_{1/2}}
\cdot t_{off}\right)
\end{equation}

\subsection{Half-Life Relationship}
\begin{equation}
t_{1/2} = \frac{\ln(2)}{\lambda}
\end{equation}

\subsection{Hierarchy Constraint}
\begin{equation}
c_{cfct} < c_{cf1t} < \min(c_{tv}, c_{pb}, c_{br})
\end{equation}

\section{Document Sources}

This synthesis draws from the following documentation files:
\begin{enumerate}
    \item \texttt{sigma\_matrix\_derivation.tex} -- Core $\Sigma = DRD$ derivation
    \item \texttt{positive\_definiteness\_constraints.tex} -- PD analysis
    \item \texttt{correlation\_structure\_design.tex} -- Block structure design
    \item \texttt{carryover\_correlation\_theory.tex} -- Carryover separation
    \item \texttt{correlation\_structure\_white\_paper.Rmd} -- AR(1) alternatives
    \item \texttt{carryover\_modeling\_white\_paper.Rmd} -- Five carryover models
    \item \texttt{mixed\_model\_comparison\_and\_best\_practices.md} -- Model guidance
    \item \texttt{technical\_differences\_scaling\_and\_pd.tex} -- Implementation
    \item \texttt{correlation\_parameters\_guide.tex} -- Parameter interpretation
    \item \texttt{whitepaper.tex} -- Crossover vs N-of-1 comparison
    \item Part XI (Alternative Summary Statistics) -- Original theoretical development
    using Fisher's $z$-transformation for closed-form power analysis
\end{enumerate}

\vspace{1cm}
\noindent\rule{\textwidth}{0.4pt}

\noindent\textit{Document generated: December 17, 2025}\\
\textit{Repository: pmsimstats2025}\\
\textit{Total mathematical documentation reviewed: 10+ documents, $\sim$15,000 lines}

\end{document}
